\clearpage
\fancyhead[L]{\small GAUSS TO VON LINDENAU.}
\hypertarget{native-306-lower}{}
\begin{center}GAUSS TO VON LINDENAU.
\end{center}

\begin{center}
\rule{4cm}{0.4pt}\\
\fontsize{10}{12}\selectfont Monatliche Correspondenz. Vol. XVIII. Pp. 83..86. July 1808.\\
\rule{4cm}{0.4pt}
\end{center}

\begin{flushright}
Göttingen, 27 June 1808.
\end{flushright}

\noindent —\;—\;— Although my work was already listed in the book fair catalogue, the printing is unfortunately still less than half finished. It is going very slowly, since I have each sh\SourceLeaf{307}eet sent to me for final revision, In connection with a problem that forms the subject of a section of this work, namely how to determine a planet’s orbit from \emph{four} geocentric positions (two of which may be incomplete), I made another calculation of the orbit of \emph{Vesta} a long time ago to give a specially chosen example, and used one of the latest Milan observations in it. Given the nature of the problem, this orbit could not and was not meant to be the most accurate result possible from the available observations (since all the positions used are \emph{individual} observations that were not first cleared of the most probable observational error). Even so, it was to be expected that these elements would be much more accurate than the III$^{\text{rd}}$, which already differed by $1'$ from the observations in September. As these new elements are still unpublished, I am sending them to you here.


\begin{center}\fontsize{9}{11}\selectfont\setlength{\tabcolsep}{3pt}\renewcommand{\arraystretch}{1.15}\begin{tabular}{lr}
Epoch, Paris meridian 1807 .\;.\;.\;.\;.\;.\;.\;.\;.\;.\;.\;.\; & $168^\circ\;10'\;47''6$\\
Daily mean tropical motion .\;.\;.\;.\;.\;.\;.\;.\;.\;.\;.\;.\;.\;.\; & $978''8588$\\
\shortstack[l]{Aphelion for the epoch\\fixed with respect to the stars} .\;.\;.\;.\;.\;.\;.\;.\; & $69^\circ\;57'\;6''5$\\
Ascending node .\;.\;.\;.\;.\;.\;.\;.\;.\;.\;.\;.\;.\;.\;.\;.\; & $103\;11\;57''3$\\
Eccentricity .\;.\;.\;.\;.\;.\;.\;.\;.\;.\;.\;.\;.\;.\;.\;.\;.\;.\;.\;.\;.\; & $0{,}0880158$\\
Logarithm of the mean distance .\;.\;.\;.\;.\;.\;.\;.\;.\;.\;.\;.\; & $0{,}3728980$\\
Inclination of the orbit .\;.\;.\;.\;.\;.\;.\;.\;.\;.\;.\;.\;.\;.\;.\;.\;.\; & $7^\circ\;8'\;20''8$\\
\end{tabular}\end{center}


\par\indent On 22 June we found Vesta again here. We have had no clear night since, but the star we observed stood almost exactly where Vesta was expected. It was of eighth magnitude (about as bright as the planet must be now), and there was no other discernible star nearby that could have been Vesta. In his earlier survey of this region, Professor \textsc{Harding} had noticed no fixed star here that had escaped observation. Finally, my own observations over about an hour already seem to show the motion in R.A. fairly unmistakably. So there can hardly be any doubt that it really was Vesta. The first clear night will give us complete certainty. We had already found and observed Juno again on 20 June, although that evening we were more inclined to identify it with another star of magnitude 10, which, however, had stayed exactly where it was on the 21$^{\text{st}}$. The observation on the 22$^{\text{nd}}$ gave us complete certainty. Juno is barely of the 10th magnitude, but I hope you will observe it with the transit instrument, since in 1806 it was observed at Milan with the mural quadrant when it was about as faint.

\Needspace{10\baselineskip}
\par\indent Here are my observations of Juno and Vesta:


\begin{center}\fontsize{9}{11}\selectfont\setlength{\tabcolsep}{3pt}\renewcommand{\arraystretch}{1.15}\begin{tabular}{lrrrr}
& 1808 & \shortstack{Mean time\\at Göttingen} & \shortstack{Apparent right\\ascension} & \shortstack{Apparent south\\declination}\\\hline
$\JunoSymbol$ & June 20 & $11^{\text{u}}\,49^{\text{m}}\,0^{\text{s}}$ & $315^\circ\;29'\;34''$ & $2^\circ\;16'\;23''$\\
$\JunoSymbol$ & 22 & $12\;0\;45$ & $315\;23\;1$ & $2\;14\;20$\\
\VestaAltarThreeOhSeven & 22 & $13\;46\;25$ & $353\;12\;17$ & $9\;11\;55$\\
\end{tabular}\end{center}


\par\indent For Juno, then, the 7$^{\text{th}}$ elements (Gött. gelehrte Anz. 1808, issue 14, Monatl. Corresp. 1808, Feb.) give a right ascension $11'$ too small and a declination $2'$ too large. For Vesta the third elements give a right ascension $8'$ too small and a declination $3'$ too large. The elements above, however, give a calculated R.A. of $353^\circ\;14'\;30''$ and a declination of $9^\circ\;12'\;26''$. The former is thus $2'\;13''$ too large and the latter $31''$ too large. The observed declination is, incidentally, rather uncertain because the comparison star was in a very unfavorable position.

I can also give you h\SourceLeaf{308}ere a few observations made this month.

\par\indent Immersion of i Virginis at the Moon’s dark limb on 4 June, $9^{\text{u}}\,3^{\text{m}}\,39^{\text{s}}0$ mean time, observed by Professor \textsc{Harding} and me in one and the same second.

\Needspace{8\baselineskip}
\par\indent Immersion of Jupiter’s II$^{\text{nd}}$ satellite on 22 June:


\begin{center}\fontsize{10}{12}\selectfont\setlength{\tabcolsep}{3pt}\renewcommand{\arraystretch}{1.15}\begin{tabular}{ll}
$12^{\text{u}}\,45^{\text{m}}\,9^{\text{s}}$ mean time & \textsc{Harding} using the $3\frac{1}{4}$-foot \textsc{Dollond}.\\
$12\;45\;32$\quad — & \textsc{Gauss} using the 17-foot \textsc{Herschel}.\\
\end{tabular}\end{center}


\par\indent We expect the fourth sheet of Professor \textsc{Harding}’s star charts soon.

